\section{The physical clock establishes the motion baseline}
The matched G33 pilot compares a 257-point relative-time representation, a native 0.25-second representation and empirical transfer of a training motion profile onto a generated path. It uses 925 training legs, two fitted seeds and one new SUMO reference episode. Of 100 declared requests, map construction and reference extraction each support 97, with 94 in their intersection. Every full pool contains eight draws per supported request. These pilot observations motivate the physical clock; they are not an equal-compute comparison or a new large-corpus temporal ablation.

Native acceleration $W_1$ falls from 0.8397/0.8319 to 0.1492/0.1575~m/s$^2$ across seeds 42/43, reductions of 82.2/81.1\%. Empirical transfer gives 0.3245/0.3588~m/s$^2$. The native model improves that endpoint in all five classes for both seeds. Native-speed distance also falls in the macro, from 0.6794/0.8238 to 0.6093/0.7437~m/s, but not in every class. One-second mean-absolute-acceleration distance decreases from approximately 0.113 to 0.0661/0.0560~m/s$^2$, providing a related observation at a coarser clock.

\fig{g33-native-motion.pdf}{1}{Matched temporal-resolution pilot: native-interval distributions for two diffusion clocks and empirical motion transfer. Both model seeds use the same request population and corresponding shared duration draws. This one-episode result motivates the clock used subsequently; it is not pooled with the final G40 comparison.}{g33-native-motion}

Longer-timescale structure remains difficult. Longest-stop distance worsens from 4.69/4.89 to 7.75/7.66 seconds; for buses it rises from 12.03/11.76 to 26.37/25.66 seconds. Lower acceleration discrepancy therefore does not establish correct stop organization. Native fitting takes 1,520/1,529 seconds compared with 301/319 seconds for the relative-time fits, at the same 10,000 updates. The complete physical and class tables, including adverse outcomes, remain in Appendix~\ref{app:native-results}.

\section{Learning route preference changes the delivered route distribution}
The G38 route experiment uses 23,696 accepted training legs and 5,748 development legs for next-edge learning. The class-equal development next-choice accuracy is approximately 96.8\% for the two neural policies, compared with 94.9\% for the shortest-cost rule. This is an observed branching-choice task; it is not an exact whole-route recovery score.

Two prospective SUMO episodes and two policy seeds compare distributions over a finite, geometrically constructed route catalog. Both episodes have 100 declared requests and 98 supported map conditions, with 97 and 92 physical/reference-route pairs. The exact realised SUMO route is in the catalog for 64 and 61 of those pairs. This coverage limitation is retained. A different permitted route between the endpoints is not automatically an erroneous prediction.

Averaged over the four cells, learned preference R reduces road-edge usage Jensen--Shannon divergence from 0.04000 for shortest-route choice to 0.03458, a 13.6\% reduction, with the same direction in all four cells. Effective route count increases from 1.000 to 1.176. Weighted native-speed and acceleration distances decrease from 1.20327/0.17435 to 1.16728/0.17192 in their respective units. These are modest distributional changes on the same public network, not evidence of arbitrary-city generalization.

\fig{g38-route-motion.pdf}{1}{Learned route preference and route/motion distribution distances in the closed G38 study. Values summarize two episodes and two seeds; the underlying class and cell results remain archived. Lower edge-use divergence is closer to the observed distribution. The plotted points show each scene/seed cell.}{g38-route-preference}

Uniform route sampling retains a relevant tradeoff: its speed distance is lower at 1.05113~m/s and its effective route count is 3.454, but edge-use divergence is higher at 0.10177. The neural route policy is therefore not a winner on every motion or diversity endpoint. Its role is a learned preference over declared feasible alternatives. Subsequent G40 arms use common R draws so that GAF effects can be attributed to conditional motion changes rather than different route selection.

\section{Both GAF views recognize classes, with different weaknesses}
The complete-cohort G35 readers use the same 29,895 training legs and 7,230 development legs as the expanded motion corpus. Spatial macro F1 is 76.47\% and speed macro F1 is 90.34\%. Accuracies are 86.17\% and 89.13\%, respectively. Table~\ref{tab:g35-five-class-readers} exposes class imbalance and the spatial reader's taxi/ride-hailing weakness; overall accuracy alone would obscure that result. These figures describe the self-trained readers before G40 joint adaptation, not the final adapted encoders.

\tab{Current GAF readers on the same development population. F1 values are percentages. Support counts are legs, not independent simulator scenes.}{g35-five-class-readers}{\begin{tabular}{lrrr}\toprule
Class & Development legs & Spatial F1 & Velocity F1 \\
Bus & 155 & 95.65 & 100.00 \\
Motorcycle & 1591 & 73.57 & 78.90 \\
Private car & 4436 & 95.20 & 92.23 \\
Bicycle & 668 & 72.10 & 99.48 \\
Taxi / ride-hailing & 380 & 45.85 & 81.07 \\
Macro F1 & 7,230 & 76.47 & 90.34 \\
\bottomrule\end{tabular}}


Spatial images classify private cars well but do not reliably separate the service-related taxi label. Speed images retain useful information about the same labels, without being added to the spatial classifier's input. Neither result proves that the two representations contain independent generative factors. The source-demand diagnosis in Chapter~\ref{ch:native-data} and the matched earlier/current-reader comparisons in the appendix provide the data context; they are not a scaling-law claim.

\section{Why stronger recognition is not sufficient feedback}
The completed G39 frozen-feature alignment study includes seven arms, two fitted seeds and eight new reference episodes. Its 800 declared episode requests yield 782 route-supported requests and 754 physical pairs; 470 realised reference routes are present in their catalogs. The dual-view SV arm gives K1 speed/acceleration distances of 1.29337/0.18761 compared with 1.28383/0.18674 for B: increases of 0.743\% and 0.466\%. The speed-only arm slightly improves primary acceleration, but the same-class shuffled control changes it by a similar amount and the four-draw weighted comparison reverses the apparent speed-only advantage. Frozen-feature supervision does not establish a general improvement of motion.

Longer-timescale changes are also mixed. SV reduces longest-stop distance by 1.44\%, while mean absolute acceleration on the one-second clock worsens by 4.86\% and four-second speed-change distance worsens by 5.23\%. Bicycle and motorcycle gains coexist with adverse bus and private-car results. The complete endpoints remain in the experiment records rather than being condensed into a favourable class score.

A separate same-pool diagnostic on the already exposed G38 scenes distinguishes global GAF weighting from weighting with fixed route and route--duration marginals. It uses four scene/seed cells and 10,784 candidates. Under the fixed joint marginals, speed-GAF weighting changes native speed/acceleration distances by only $-0.12\%/-0.13\%$, with improvements in one and two of the four cells respectively. Spatial and dual-view weighting worsen both macro endpoints. The preserved marginals are distributional constraints; an individual coupled random draw need not retain the same route and duration.

\fig{g39-fixed-marginals.pdf}{1}{Closed mechanism diagnostic: changes under global, fixed-route and fixed-route--duration weighting of identical candidate pools. The last comparison preserves the joint route/time marginal rather than asserting that each selected sample is identical. These reused scenes are development evidence; the small residual speed-view effects are not independent confirmation of guidance.}{g39-fixed-marginals}

Independent augmentation nevertheless gives a useful, narrower finding. Across two raw-series learner seeds, real-only training reaches 91.25\% macro F1, real plus baseline diffusion data 93.17\%, real plus speed-aligned data 94.10\%, and real plus RAW-aligned data 94.06\%. Real plus dual-view data reaches 93.57\%; synthetic-only regimes remain below real-only. Thus diffusion samples can help this downstream task, while the almost equal speed-GAF and RAW results do not demonstrate an image-specific advantage. The generator has access to the full labelled corpus, so this is not a strict low-label learning claim.

\fig{g39-alignment-utility.pdf}{1}{Completed G39 frozen-feature alignment (left) and independent augmentation (right), with both downstream seeds shown. The evaluation model reads raw time series and is independent of the GAF readers. Better task performance does not override adverse physical endpoints or establish that GAF is necessary.}{g39-alignment-utility}

These studies motivate the final question: can a GAF encoder trained as part of a usable generative condition perform better than a fixed recognition score or an auxiliary frozen feature target? They do not predetermine the answer.
